\section{Introduction}
We set $q\in\mathbb{C}^*$ and $0<|q|<1$ throughout this paper. Basic hypergeometric series with the base $q$ is defined by
\begin{gather}\label{qHGS}
{}_r\varphi_s\left(\bm{a};\bm{b};q,x\right)=
{}_r\varphi_s\left(\begin{matrix}\bm{a}\\ \bm{b}\end{matrix};q,x\right)
:=\sum_{n\ge0} \frac{(a_1,a_2,\ldots,a_r;q)_n}{(b_1,\ldots,b_s;q)_n(q;q)_n}\big\{(-1)^nq^{\frac{n(n-1)}{2}}\big\}^{1+s-r}x^n,
\end{gather}
where $x\in\mathbb{C}$, $\bm{a}=(a_1,a_2,\ldots,a_r)\in\mathbb{C}^r$, $\bm{b}=(b_1,\ldots,b_s)\in\mathbb{C}^s$ and $b_1,\ldots,b_s\notin q^{-\mathbb{N}}$.

Basic hypergeometric series \eqref{qHGS} is a $q$-analog of the classical hypergeometric series
\begin{gather}\label{HGS}
{}_rF_s\left(\bm{\alpha};\bm{\beta};x\right)=
{}_rF_s\left(\begin{matrix}
\bm{\alpha}\\
\bm{\beta}
\end{matrix};x\right)
:=\sum_{n\ge0}\frac{(\alpha_1)_n(\alpha_2)_n\cdots(\alpha_r)_n}{(\beta_1)_n\cdots(\beta_s)_nn!}x^n,
\end{gather}
where $\bm{\alpha}=(\alpha_1,\alpha_2,\ldots,\alpha_r)\in\mathbb{C}^r$, $\bm{\beta}=(\beta_1,\ldots,\beta_s)\in\mathbb{C}^s$ and $\beta_1,\ldots,\beta_s\notin \mathbb{Z}_{\le0}$.

The radius of convergence of the series \eqref{qHGS} and \eqref{HGS} are both $\infty$, $1$ or $0$ according to whether $r<s+1$, $r=s+1$ or $r>s+1$. In this paper, we assume $r>s+1$ and $a_1a_2\cdots a_rb_1b_2\cdots b_s\neq0$. Therefore the series \eqref{qHGS} and \eqref{HGS} are divergent in this paper.

Basic hypergeometric series \eqref{qHGS} formally satisfies the following linear $q$-difference equation:
\begin{gather}\label{qHGEq}
\left[x(-\sigma_q)^{1+s-r}\prod_{j=1}^r(1-a_j\sigma_q)-(1-\sigma_q)\prod_{k=1}^s\left(1-\frac{b_k}{q}\sigma_q\right)\right]y(x)=0,
\end{gather}
where $\sigma_q$ is the $q$-shift operator defined by $\sigma_qy(x)=y(qx)$. As applying the local theory of linear $q$-difference equations (cf.\ Adams~\cite{Adams1929}), we see that the equation~\eqref{qHGEq} has a fundamental system of solutions around infinity\footnote{In his paper, Adams used the function $q^{-\frac{1}{2}(\log_q x)(\log_q x-1)}$ instead of $\theta_q(x)$. We remark that two functions $q^{-\frac{1}{2}(\log_q x)(\log_q x-1)}$ and $\theta_q(x)$ play the same role in constructing formal solutions since they satisfy the same $q$-difference equation
\begin{gather*}
\sigma_q^n y=x^{-n}q^{-\frac{n(n-1)}{2}}y.
\end{gather*}}:
\begin{gather}\label{qHGEq_Sol}
y_i^{(\infty)}(x)=\frac{\theta_q(-a_ix)}{\theta_q(-x)}{}_{r}\varphi_{r-1}
\left(\begin{matrix}
a_i,\frac{a_iq}{b_1},\frac{a_iq}{b_2},\ldots,\frac{a_iq}{b_s},0,\ldots,0\\
\frac{a_iq}{a_1},\ldots,\frac{a_iq}{a_{i-1}},\frac{a_iq}{a_{i+1}},\ldots,\frac{a_iq}{a_r}
\end{matrix};q,\frac{qb_1\cdots b_s}{a_1\cdots a_rx}\right),\qquad\!\! 1\le i\le r,\!\!\!
\end{gather}
which are convergent in $\vert x\vert>\vert qb_1\cdots b_s/a_1\cdots a_r\vert$.

In this paper, we construct an actual solution of \eqref{qHGEq} which admits \eqref{qHGS} as an asymptotic expansion and give its analytic continuation, by using $q$-Borel--Laplace transform (see Theorem~\ref{Theorem:Main1}). This result shows a $q$-analog of Stokes phenomenon. After that, we consider to take the limit $q\to1$ in Theorem~\ref{Theorem:Main1} (see Theorem~\ref{Theorem:Main2}).

The motivation of this study comes from as follows: Ichinobe \cite{Ichinobe2001} obtained the Borel sum of divergent classical hypergeometric series.
\begin{Theorem}[Ichinobe {\cite[Theorem 2.1]{Ichinobe2001}}]\label{Theorem:Ichinobe}
Let $\bm{\alpha}=(\alpha_1,\alpha_2,\ldots,\alpha_r)\in\mathbb{C}^r$ and $\bm{\beta}=(\beta_1,\beta_2,\ldots,$ $\beta_s)\in\mathbb{C}^s$. Assume $\alpha_i-\alpha_j\notin\mathbb{Z}$, $i\neq j$. Then ${}_rF_s(\bm{\alpha};\bm{\beta};x)$ is $1/(r-s-1)$-summable in any direction $d$ such that $d\neq0$ $({\rm mod}~2\pi)$ and its Borel sum $f(x)$ is given by
\begin{gather*}\label{Ichinobe_result1}
f(x)=C_{\bm{\alpha\beta}}\sum_{j=1}^rC_{\bm{\alpha\beta}}(j)(-x)^{-\alpha_j}{}_{s+1}F_{r-1}\left( \begin{matrix}
\alpha_j,1+\alpha_j-\bm{\beta}\\
1+\alpha_j-\widehat{\bm{\alpha}_j}
\end{matrix}
;\frac{(-1)^{1+s-r}}{x}\right),
\end{gather*}
where $x\in S(\pi,(r-s+1)\pi,\infty):=\{x\in\mathbb{C}^*;\left|\pi-\arg x\right|<(r-s+1)\pi/2)\}$ and $\widehat{\bm{\alpha}_j}\in\mathbb{C}^{r-1}$ is the vector which is obtained by omitting the $j$-th component from $\bm{\alpha}$ and
\begin{gather}\label{Ichinobe_result2}
C_{\bm{\alpha\beta}}=\frac{\Gamma(\bm{\beta})}{\Gamma(\bm{\alpha})},\qquad C_{\bm{\alpha\beta}}(j)=\frac{\Gamma(\alpha_j)\Gamma(\widehat{\bm{\alpha}_j}-\alpha_j)}{\Gamma(\bm{\beta}-\alpha_j)}.
\end{gather}
Here we use the following abbreviations
\begin{gather*}
\Gamma(\bm{\alpha})=\prod_{l=1}^r\Gamma(\alpha_l),\qquad \Gamma(\widehat{\bm{\alpha}_j}-\alpha_j)=\prod_{l=1,\,l\neq j}^r\Gamma(\alpha_l-\alpha_j).
\end{gather*}
\end{Theorem}
Since basic hypergeometric series \eqref{qHGS} is a $q$-analog of hypergeometric series \eqref{HGS}, whether there exists a $q$-analog of Theorem \ref{Theorem:Ichinobe} is a natural question. To answer this question, we use the theory of $q$-Borel summability. It is a $q$-analog of the theory of Borel summability and studied by many authors (cf.~Ramis \cite{Ramis1992}, Zhang \cite{Zhang2002,Zhang2005}, Di Vizio--Zhang \cite{Vizio-Zhang} and Dreyfus--Eloy~\cite{DreyfusEloy2016}). Dreyfus~\cite{Dreyfus2015_1} proved that for every formal power series solution of a linear $q$-difference equation with rational coefficients, we may construct a meromorphic solution of the same equation applying several $q$-Borel and $q$-Laplace transformations of appropriate orders and appropriate direction.

Our results are generalizations of Zhang \cite{Zhang2002} and Morita \cite{Morita2014}. Zhang studied the $q$-Borel summability of divergent basic hypergeometric series ${}_2\varphi_0$ which is the case of $r=2$ and $s=0$ in our result. Later, Morita obtained similar results with Zhang for ${}_3\varphi_1$, which is the case of $r=3$ and $s=1$ in our result. The common point of their assumptions is that $r$ and $s$ satisfy $r-s=2$. Our result gives the resummation of ${}_r\varphi_s$ in the case of $r-s\ge2$.

This paper is organized as follows. We fix our notions and review the theory of $q$-Borel summability in Section \ref{Section:Preliminary}. Main results of this paper are given in Section \ref{Section:Main Results}. In Sections~\ref{Section:Proof_Main1} and~\ref{Section:Proof_Lemma:Main1}, we give proofs of Theorem~\ref{Theorem:Main1} and Lemma~\ref{Lemma:Main1} respectively. A~proof of Theorem~\ref{Theorem:Main2} is given in Section~\ref{Section:Proof_Main2}.

